\subsection{Clause Resolution}

Although we've already seen the pseudocode for clause resolution, it's not discussed very 
much, and I want us to understand how it works better. As we make decisions and propagations, 
we form this thing called an \textbf{implication graph}. In a DPLL style solver, our assignments 
and propagations form a directed acyclic graph (I will not be proving this), and by traversing the
implication graph, we can find out what decisions led to the contradiction.

To learn more about implication graphs, let's look at one and discuss what it means.
We will use the implication graphs from here (Figures 4.1 and 4.2): 
\begin{center}
    \small{\url{https://www3.cs.stonybrook.edu/~cram/cse505/Fall20/Resources/cdcl.pdf}}
\end{center}

We will start with Figure 4.1:

\fbox{
    \begin{tikzpicture}[
        >=Stealth,
        every node/.style={font=\large},
        implication/.style={->, thick}
    ]

    % Nodes
    \node (x31) {$x_{31}=0\ @\ 3$};
    \node (x2)  [below right=0.8cm and 1.8cm of x31] {$x_2=0\ @\ 5$};
    \node (x1)  [below left=0.8cm and 1.8cm of x2] {$x_1=0\ @\ 5$};
    \node (x3)  [below right=0.8cm and 1.8cm of x1] {$x_3=0\ @\ 5$};
    \node (x4)  [above right=0.8cm and 1.8cm of x3] {$x_4=1\ @\ 5$};

    % Edges
    \draw[implication] (x31) -- node[above left] {$\omega_1$} (x2);
    \draw[implication] (x1)  -- node[above left] {$\omega_1$} (x2);
    \draw[implication] (x1)  -- node[below left] {$\omega_2$} (x3);
    \draw[implication] (x2)  -- node[above right] {$\omega_3$} (x4);
    \draw[implication] (x3)  -- node[below right] {$\omega_3$} (x4);

    \end{tikzpicture}
}

It corresponds to the following CNF equation: 
$$ (\omega_1) \wedge (\omega_2) \wedge (\omega_3) =
(x_1 \vee x_{31} \vee \neg x_2) \wedge (x_1 \vee \neg x_3) \wedge (x_2 \vee x_3 \vee x_4)$$

Each node of the graph follows this form: $x_i = v @ d$, but what does it mean?

\begin{center}
    This is the notation for assignment of a variable at a certain decision level:
    $$x_i = v@d$$
    where $i$ denotes the index of the variable being assigned, $v \in \{0, 1\}$
    (alternatively, $v \in \{F, T\}$), and $d$ is the decision level for which that 
    assignment was executed.
\end{center}

So in our graph, $x_{31} = 0 @ 3$ means variable 31 was assigned 0 (or False) at
decision level 3. What about the edges? The edges are clauses, and they connect
variables to each other, but what's the relationship? Let's go left to right in our 
implication graph. Variables 31 and 1 are assigned False at decision levels 3 and 5
respectively, and through clause 1, they both have an edge to variable 2, which is also 
assigned False at decision level 5. If we look at the original CNF equation, we see that 
in clause 1, $x_1$ and $x_{31}$ are not negated, so they are False literals. That leaves 
$x_2$ Unassigned, so we propagate in order to satisfy the clause. We assign $x_2 = F$ 
because it's negated, so its literal will evaluate to True. Likewise, $x_1$ has another 
edge of clause 2 ($\omega_2$) to $x_3$. If we look at clause 2, we see that $x_1 = F$ 
forces $\neg x_3 = T \equiv x_3 = F$. Lastly, we see that $x_2$ and $x_3$ have an edge 
to $x_4$ through $\omega_3$, forcing its assignment to be $x_4 = T$.

That was pretty easy to understand, so now let's look at an implication graph with 
contradictions. From Figure 4.2:

\noindent
\fbox{
\resizebox{\textwidth}{!}{%
\begin{tikzpicture}[
    >=Stealth,
    every node/.style={font=\large},
    implication/.style={->, thick}
]

% Left side
\node (x31) {$x_{31}=0\ @\ 3$};

\node (x2)
    [below right=0.8cm and 1.8cm of x31]
    {$x_2=0\ @\ 5$};

\node (x1)
    [below left=0.8cm and 1.8cm of x2]
    {$x_1=0\ @\ 5$};

\node (x3)
    [below right=0.8cm and 1.8cm of x1]
    {$x_3=0\ @\ 5$};

\node (x4)
    [above right=0.8cm and 1.8cm of x3]
    {$x_4=1\ @\ 5$};

% Right side
\node (x5)
    [above right=0.8cm and 2.5cm of x4]
    {$x_5=0\ @\ 5$};

\node (x6)
    [below right=0.8cm and 2.5cm of x4]
    {$x_6=0\ @\ 5$};

\node (x21)
    [below left=1.0cm and 1.8cm of x6]
    {$x_{21}=0\ @\ 2$};

\node (kappa)
    [right=2.5cm of x6]
    {$\kappa$};

% Left implication graph
\draw[implication]
    (x31) -- node[above left] {$\omega_1$} (x2);

\draw[implication]
    (x1) -- node[above left] {$\omega_1$} (x2);

\draw[implication]
    (x1) -- node[below left] {$\omega_2$} (x3);

\draw[implication]
    (x2) -- node[above right] {$\omega_3$} (x4);

\draw[implication]
    (x3) -- node[below right] {$\omega_3$} (x4);

% Right implication graph
\draw[implication]
    (x4) -- node[above left] {$\omega_4$} (x5);

\draw[implication]
    (x4) -- node[below left] {$\omega_5$} (x6);

\draw[implication]
    (x21) -- node[below right] {$\omega_5$} (x6);

\draw[implication]
    (x5) -- node[above right] {$\omega_6$} (kappa);

\draw[implication]
    (x6) -- node[below right] {$\omega_6$} (kappa);

\end{tikzpicture}
}
}

\noindent
This graph corresponds to the CNF equation: 
$$ (\omega_1) \wedge (\omega_2) \wedge (\omega_3) \wedge (\omega_4) \wedge (\omega_5) \wedge (\omega_6) 
=
(x_1 \vee x_{31} \vee \neg x_2)_1 \wedge (x_1 \vee \neg x_3)_2 \wedge (x_2 \vee x_3 \vee x_4)_3
$$
$$
\wedge (\neg x_4 \vee \neg x_5)_4 \wedge (x_{21} \vee \neg x_4 \vee \neg x_6)_5 \wedge (x_5 \vee x_6)_6
$$

Up until $x_4 = 1 @ 5$, this graph is the same as the one from above. After that,
we have some more decisions and propagations, that I bet you can read by now, up until 
we reach this symbol: $\kappa$, which has two edges via $\omega_6$. This node, $\kappa$,
corresponds to a contradiction within clause $\omega_6$ because its literals are all False.

\noindent
When we resolve this conflict, we first start with the clause 
$$(x_5 \vee x_6)$$
as that is the clause in which our conflict occurred. Next, we find the most recently propagated literal, 
which we will assume to be $x_6$ (the most recently propagated literal depends on our implementation 
of propagation, so for simplicity's sake, we will assume numerical order to settle ties in 
implication orderings on the same level). $x_6$ was propagated because of the clause
$$(x_{21} \vee \neg x_4 \vee \neg x_6)$$
Combining this with the current conflict clause yields:
$$(x_5 \vee x_6 \vee x_{21} \vee \neg x_4 \vee \neg x_6)$$
which we process to remove all instances of the conflicting variable, leading to a resolution 
clause of: 
$$(x_5 \vee x_{21} \vee \neg x_4)$$
In this clause, the most recently propagated literal is $x_5$, so we grab its reason clause ($\omega_4$)
and combine once again to get: 
$$(\neg x_4 \vee \neg x_5 \vee x_5 \vee x_{21} \vee \neg x_4)$$

\noindent
Simplifying that clause gives us: $$(\neg x_4 \vee x_{21})$$
The only propagated literal in the 
clause is $x_4$, so we grab its reason clause ($\omega_3$) and combine again into:
$$(x_2 \vee x_3 \vee x_4 \vee \neg x_4 \vee x_{21})$$
This simplifies down into:
$$(x_2 \vee x_3 \vee x_{21})$$
Our most recently propagated literal is now $x_3$. Repeating 
the process again, we get an unsimplified resolution clause of: 
$$(x_2 \vee x_3 \vee x_{21} \vee x_1 \vee \neg x_3)$$ 
Simplifying it down, we now get:
$$(x_2 \vee x_{21} \vee x_1)$$
Our last propagated literal in the clause is $x_2$. We 
combine our incomplete resolution clause with $\omega_1$ to get:
$$(x_2 \vee x_{21} \vee x_1 \vee x_1 \vee x_{31} \vee \neg x_2)$$
This transforms into 
$$(x_1 \vee x_{21} \vee x_{31})$$
which contains only decision literals.

What does this new clause, which we'll call $\omega_7$, tell us? It tells us that 
all three decision literals \textbf{cannot} all be False (ie, at least one of them 
\textbf{must} be True).

\subsection{Clause Learning - Results}

Here are the results for Clause Learning:

\noindent
uf20.91: \textbf{425ms}, up 98ms from Iterative DPLL 

\noindent
UUF50.218.1000: \textbf{9652ms}, up 3287ms from Iterative DPLL

Why did our runtime go up? Well, clause learning isn't easy, and the more 
contradictions we encounter, the more expensive clause learning we're performing. 
Another reason is that we're increasing the memory load every time we add in a new
clause. Luckily, our testing sets are small, but on large problems, this would be a 
much larger issue. So then why do modern solvers use clause learning? Because there 
are ways to make it much faster. The first step would be to implement non-chronological 
backjumping, where instead of moving up just one decision level in our trail, we move up 
multiple at a time. The second step would be to implement First UIPs, which makes clause 
learning a less expensive operation (doesn't fix the memory load though), and it
combined with non chronological backjumping allows us to instantly and fully use our 
learned clause to its potential. Although we have taken a small step back, we will 
see what leaps forward we will make in the next sections.

\textit{Aside: Why did we get rid of pure literal elimination? The problem with pure literal 
elimination is that once we start learning clauses, previously pure literals may no longer 
be pure anymore once we add in a new clause. Because pure literals aren't born from decisions/
assignments, they aren't as ``concrete'' in their status. (Ie, propagation \textbf{forces} new 
literal assignments, pure literal elimination is more of a ``rule of thumb''). If we had a pure 
literal, say $5$, and we added in a clause containing $-5$, unfortunately, $5$ is no longer pure, 
and anything propagated because of $5$ also would need to be undone because the reason why they 
were propagated/assigned those values in the first place doesn't hold anymore. Because we saw that 
pure literal elimination wasn't that helpful to begin with and with all these reasons as to why it 
wouldn't easily work under CDCL, it would just be much better to not have it at all.}

\begin{FVerbatim}
On all 1000 equations in uf20-91:

Clause Learning:           425ms (~0.43 seconds).
Iterative DPLL (AbP):      328ms (~0.33 seconds).
Iterative DPLL (PbA):      324ms (~0.32 seconds).
Watched Literals:          349ms (~0.35 seconds).
Pure Literal Elimination:  1548ms (~1.55 seconds).
Unit Propagation:          868ms (~0.87 seconds).
Recursive Backtracking:    1,540,613ms (~25.68 minutes).
Brute Force:               1,770,851ms (~29.51 minutes).
Short-Circuited (Basic):   49,874ms (~50 seconds).
Parallel:                  10,589ms (~10.5 seconds).

On all 1000 equations in UUF50.218.1000:

Clause Learning:          9,652ms (~9.65 seconds).
Iterative DPLL (AbP):     6,369ms (~6.37 seconds).
Iterative DPLL (PbA):     6,365ms (~6.37 seconds).
Watched Literals:         7,130ms (~7.13 seconds).
Pure Literal Elimination: 231,107ms (~3.85 minutes).
Unit Propagation:         140,041ms (~2.34 minutes).
\end{FVerbatim}
